\begin{abstract}
\noindent
La boiterie est l'un des trois problèmes de santé les plus fréquents chez les vaches
laitières, avec des prévalences de 20 à 55\,\% selon les études. Son coût économique
(75 à 300+ USD par cas) et son impact sur le bien-être animal motivent le développement
de systèmes de détection automatisée. Ce mémoire présente la conception,
l'implémentation et la validation d'un système reproductible de détection de boiterie
probable, à partir de données de capteurs comportementaux collectées sur 28~vaches
pendant 41~jours, comme l'activité, les transitions posturales et la dynamique de
mouvement. L'architecture proposée est modulaire~: un \textit{pipeline} en quatre étapes
(chargement, extraction d'environ 50~variables dérivées, détection d'anomalies,
application de règles métier) combine un détecteur d'anomalies non supervisé proposé
(Isolation Forest) avec des règles métier explicites comme la persistance temporelle,
la cohérence multi-signaux et l'espacement des notifications. L'implémentation comprend
également un score continu de risque pour le classement des animaux et un tableau de
bord interactif offrant à l'éleveur quatre vues complémentaires.

La validation du système proposé est conduite sur deux niveaux. Une validation manuelle
par injections contrôlées confirme que les événements détectables sont récupérés avec
un recouvrement temporel strict, tandis que les scénarios non détectables restent
silencieux. Une validation robuste multi-configurations, avec plus de 9\,500~exécutions et
336~paramétrisations, établit 100\,\% de détection sur les cas forts, une fuite nulle
et un taux de fausses notifications limité à 0.105 par vache-jour. De plus, l'accord
de décision entre deux campagnes successives atteint 91.7\,\%.

D'autre part, pour valider les choix d'architecture, une étude d'ablation portant sur
90~cas par variante compare quatre configurations du pipeline et montre que chaque
composante est nécessaire~: le détecteur seul manque de précision temporelle, tandis
que les règles seules génèrent significativement plus de fausses alertes. Isolation
Forest et Local Outlier Factor présentent en outre des profils complémentaires,
ouvrant la perspective d'un ensemble hybride. Un \textit{benchmark} de performance confirme
enfin une scalabilité linéaire ($\approx$0.6~seconde par vache).

La rigueur statistique est renforcée par des tests appariés (McNemar exact et Wilcoxon
des rangs signés), des intervalles de confiance à 95\,\% par rééchantillonnage
\textit{bootstrap}, des courbes de sensibilité/spécificité et de précision/rappel,
une analyse de l'importance des variables par valeurs de Shapley, ainsi qu'une
calibration isotonique des scores (score de Brier = 0.087~; erreur de calibration
espérée = 0.015). Une matrice de responsabilités documente l'intégration du système
dans le \textit{workflow} de ferme.

Une validation exploratoire sur un corpus complémentaire de 30~vaches (même ferme,
capteur IceTag, automne~2019) confirme la portabilité technique du pipeline et une
cohérence partielle avec les scores cliniques locomoteurs disponibles pour 12~d'entre
elles. L'ensemble constitue un cadre reproductible pour la détection non supervisée de
boiterie probable, avec traçabilité par empreintes cryptographiques des artefacts.

\vspace{1em}
\textbf{Mots clés:} boiterie bovine, données capteurs, détection d'anomalies,
Isolation Forest, règles métier, validation robuste, reproductibilité, élevage de
précision.
\end{abstract}

\begin{otherlanguage}{english}
\renewcommand{\abstractname}{Abstract}
\begin{abstract}
Lameness is one of the three most prevalent health issues in dairy cattle, affecting
20--55\,\% of herds depending on the study. Its economic cost (75--300+ USD per case)
and welfare impact motivate the development of automated detection systems. This thesis
presents the design, implementation and validation of a reproducible system for probable
bovine lameness detection, using behavioral sensor data collected from 28~cows over
41~days, such as activity, postural transitions and movement dynamics. The proposed
architecture is modular~: a four-stage pipeline (loading, extraction of about
50~derived features, anomaly detection, application of business rules) combines a
proposed unsupervised anomaly detector (Isolation Forest) with explicit business rules
such as temporal persistence, multi-signal coherence and notification spacing. The
implementation also includes a continuous risk score for ranking animals and an
interactive dashboard offering four complementary views to the farmer.

Validation of the proposed system is conducted at two levels. A manual validation
using controlled injections confirms that detectable events are recovered with strict
temporal overlap, while undetectable scenarios remain silent. A robust
multi-configuration validation, with more than 9,500~runs and
336~parameterizations, establishes 100\,\% detection on strong cases, zero leakage
and a false notification rate limited to 0.105 per cow-day. Furthermore, decision
agreement between two successive campaigns reaches 91.7\,\%.

To validate the architectural choices, an ablation study of 90~cases per variant
compares four pipeline configurations and shows that each component is necessary~:
the detector alone lacks temporal precision, whereas rules alone generate significantly
more false alerts. Isolation Forest and Local Outlier Factor further exhibit
complementary profiles, opening the prospect of a hybrid ensemble. A performance
benchmark finally confirms linear scalability ($\approx$0.6~second per cow).

Statistical rigor is reinforced by paired tests (exact McNemar and Wilcoxon
signed-rank), 95\,\% bootstrap confidence intervals, sensitivity/specificity and
precision/recall curves, a Shapley-value feature importance analysis, and isotonic
score calibration (Brier score = 0.087~; expected calibration error = 0.015). A
responsibility matrix also documents the integration of the system into the farm
workflow.

An exploratory validation on a complementary cohort of 30~cows (same farm, IceTag
sensor, fall~2019) confirms the technical portability of the pipeline and partial
concordance with locomotor clinical scores available for 12~cows. The work
constitutes a reproducible framework for unsupervised detection of probable lameness,
with cryptographic-hash traceability of all artifacts.

\vspace{1em}
\textbf{Keywords:} bovine lameness, sensor data, anomaly detection, Isolation Forest,
business rules, robust validation, reproducibility, precision livestock farming.
\end{abstract}
\end{otherlanguage}
